\chapter{Bootstrap, jackknife i testy permutacyjne: wnioskowanie przez ponowne próbkowanie}
\label{ch:resampling}
\paragraph{Powiązania i odsyłacze.} Porównaj estymatory z rozdziału~\ref{ch:estimators}, metody asymptotyczne z rozdziału~\ref{ch:asymptotic-estimation} i analizę procesów zależnych w rozdziale~\ref{ch:time-series}.


\section{Dlaczego ponownie losujemy dane?}
W wielu problemach znamy statystykę $T_n=T(X_1,\ldots,X_n)$,
ale nie znamy jej rozkładu lub wyprowadzenie go wymaga trudnych
rachunków. \emph{Resampling} polega na odtwarzaniu rozkładu
statystyki przez wielokrotne konstruowanie sztucznych prób.
Najpierw trzeba jednak ustalić, \emph{co jest jednostką losowania}.
Gdy jednostkami są pacjenci, a każdy pacjent ma dziesięć pomiarów,
nie wolno bez uzasadnienia losować niezależnie wszystkich pomiarów:
niszczy to zależność wewnątrz pacjenta.

W tym rozdziale \emph{liczebność rzeczywistej próby} oznaczamy przez $n$,
a \emph{liczbę replik resamplingowych} przez $B$. Zwiększenie $B$
zmniejsza błąd symulacyjny wyznaczania rozkładu resamplingowego,
ale nie zwiększa informacji statystycznej zawartej w $n$ obserwacjach.

\section{Bootstrap nieparametryczny: definicja i dokładny rachunek}
\begin{definicja}[Dystrybuanta empiryczna i próba bootstrapowa]
\label{def:bootstrap-sample}
Niech $(X_i)_{i\ge1}$ będą niezależnymi zmiennymi losowymi
o jednakowym rozkładzie $P$ na $(\R,\mathcal B(\R))$.
Po zaobserwowaniu $x_1,\ldots,x_n$ definiujemy
$P_n=n^{-1}\sum_{i=1}^n\delta_{x_i}$.
\emph{Próbka bootstrapowa} $(X_1^*,\ldots,X_n^*)$ to,
warunkowo względem danych, $n$ niezależnych losowań z $P_n$.
Piszemy $\Prob^*(\cdot)=\Prob(\cdot\mid X_1,\ldots,X_n)$,
$\E^*$ i $\Var^*$ odpowiednio dla warunkowej wartości
oczekiwanej i wariancji, a $\mathcal L^*$ dla warunkowego
rozkładu (względem ustalonych danych). Gwiazdką oznaczamy statystykę
wyznaczoną z próby bootstrapowej.
\end{definicja}

\begin{twierdzenie}[Dokładna wariancja średniej bootstrapowej]
\label{thm:bootstrap-mean-variance}
Niech $n\ge2$, $\E X_1^2<\infty$,
$\overline X_n=n^{-1}\sum_{i=1}^nX_i$, a
$S_n^2=(n-1)^{-1}\sum_{i=1}^n(X_i-\overline X_n)^2$.
Wówczas, dla niemal wszystkich realizacji danych,
\[
 \E^*\overline X_n^*=\overline X_n,\qquad
 \Var^*(\overline X_n^*)=\frac{n-1}{n^2}S_n^2.
\]
W szczególności bootstrapowy błąd standardowy $\sqrt{\Var^*(\overline X_n^*)}$
nie jest przy skończonym $n$ identyczny z klasycznym $S_n/\sqrt n$,
ale ich iloraz, przy $S_n>0$, wynosi $\sqrt{(n-1)/n}$.
\end{twierdzenie}
\begin{proof}
Warunkowo względem danych każda $X_i^*$ ma średnią
$\overline X_n$ i wariancję
$n^{-1}\sum_{j=1}^n(X_j-\overline X_n)^2=(n-1)S_n^2/n$.
Warunkowa niezależność $n$ losowań daje wskazaną wariancję średniej.
\end{proof}

\begin{uwaga}[Parametryczny i nieparametryczny bootstrap]
W wariancie \emph{nieparametrycznym} losujemy z rozkładu empirycznego $P_n$.
W wariancie \emph{parametrycznym} najpierw dopasowujemy rodzinę
$\{P_\theta:\theta\in\Theta\}$, a następnie losujemy z
$P_{\widehat\theta_n}$. Pierwszy wariant nie wymaga wyboru rodziny,
lecz może być mało wiarygodny dla ekstremalnych kwantyli,
gdy w danych nie ma obserwacji dalekiego ogona. Drugi wymaga
wiarygodnego modelu; symulowanie z rozkładu błędnie dobranego
może prowadzić do błędnych przedziałów mimo ogromnego $B$.
\end{uwaga}

\begin{twierdzenie}[Warunkowe centralne twierdzenie graniczne dla bootstrapu średniej]
\label{thm:bootstrap-clt}
Niech $(X_i)_{i\ge1}$ będą niezależne o tym samym rozkładzie,
$\E|X_1|^3<\infty$ i $\sigma^2=\Var(X_1)>0$.
Dla prawie każdej nieskończonej realizacji $(X_i)_{i\ge1}$ zachodzi
\[
 \mathcal L^*\bigl(\sqrt n(\overline X_n^*-\overline X_n)\bigr)
 \Rightarrow\Normal(0,\sigma^2),\qquad n\to\infty.
\]
Jest to ten sam graniczny rozkład co $\sqrt n(\overline X_n-\E X_1)$.
\end{twierdzenie}
\begin{proof}
Oznaczmy $Y_{n,j}=X_j-\overline X_n$ i
$m_{k,n}=n^{-1}\sum_{j=1}^n|Y_{n,j}|^k$ dla $k=2,3$.
Z mocnego prawa wielkich liczb $m_{2,n}\to\sigma^2$
oraz $m_{3,n}=O(1)$ prawie na pewno.
Warunkowa funkcja charakterystyczna jednej wycentrowanej
obserwacji bootstrapowej, podzielonej przez $\sqrt n$, wynosi
\[
 \psi_n(t)=\frac1n\sum_{j=1}^n
 e^{itY_{n,j}/\sqrt n}
 =1-\frac{t^2m_{2,n}}{2n}+R_n(t),\qquad
 |R_n(t)|\le \frac{|t|^3m_{3,n}}{6n^{3/2}},
\]
gdyż $\sum_jY_{n,j}=0$. Dla ustalonego $t$ i niemal każdej
realizacji $n\log\psi_n(t)\to-t^2\sigma^2/2$.
Warunkowa niezależność daje funkcję charakterystyczną sumy
równą $\psi_n(t)^n$, więc dąży ona do
$\exp(-t^2\sigma^2/2)$. Z twierdzenia Lévy'ego
o ciągłości funkcji charakterystycznych wynika teza.
Zwykłe CLT dla średniej z próby stosujemy oddzielnie.
\end{proof}

\begin{wniosek}[Metoda delty dla gładkich funkcji średniej]
\label{cor:bootstrap-delta}
W założeniach twierdzenia~\ref{thm:bootstrap-clt} niech
$g\colon\R\to\R$ będzie różniczkowalna w
$\mu=\E X_1$, przy czym $g'$ jest ciągła w otoczeniu $\mu$.
Wówczas dla prawie każdej realizacji danych
\[
 \mathcal L^*\bigl(\sqrt n[g(\overline X_n^*)-g(\overline X_n)]\bigr)
 \Rightarrow \Normal(0,\sigma^2g'(\mu)^2).
\]
Jeżeli $g'(\mu)=0$, granica jest zdegenerowana i nie należy
bez dalszej analizy stosować zwykłego przybliżenia normalnego.
\end{wniosek}
\begin{proof}
Mocne prawo wielkich liczb daje $\overline X_n\to\mu$,
a twierdzenie~\ref{thm:bootstrap-clt} zapewnia warunkową
ciasność $\sqrt n(\overline X_n^*-\overline X_n)$.
Z twierdzenia o wartości średniej i ciągłości $g'$ otrzymujemy
\[
 g(\overline X_n^*)-g(\overline X_n)
 =g'(\mu)(\overline X_n^*-\overline X_n)
   +o_{\Prob^*}(n^{-1/2})
\]
dla niemal każdej ustalonej realizacji danych.
Wystarczy zastosować warunkowe twierdzenie Slutsky'ego.
\end{proof}

\section{Przedziały bootstrapowe: trzy różne konstrukcje}
Niech $\widehat\theta_n=T_n(X_1,\ldots,X_n)$,
a $\widehat\theta_n^{*(1)},\ldots,\widehat\theta_n^{*(B)}$
będą statystykami z $B$ niezależnych prób bootstrapowych,
warunkowo względem danych. Przez $q_\gamma^*$ oznaczamy
warunkowy kwantyl rzędu $\gamma$ rozkładu
$\widehat\theta_n^*$, w praktyce szacowany kwantylem
z $B$ replik.

\begin{definicja}[Przedział percentylowy i podstawowy]
\label{def:bootstrap-intervals}
Dla $0<\alpha<1$ definiujemy \emph{przedział percentylowy}
$I_{\rm perc}=[q_{\alpha/2}^*,q_{1-\alpha/2}^*]$
oraz \emph{przedział podstawowy} (basic)
\[
 I_{\rm basic}=
 [2\widehat\theta_n-q_{1-\alpha/2}^*,
  2\widehat\theta_n-q_{\alpha/2}^*].
\]
Jeżeli istnieje dodatnie oszacowanie $\widehat s_n$ błędu
standardowego i analogiczne $\widehat s_n^*$ w próbie bootstrapowej,
można też użyć \emph{przedziału studentyzowanego},
opartego na kwantylach
$(\widehat\theta_n^*-\widehat\theta_n)/\widehat s_n^*$.
Wszystkie te konstrukcje wymagają osobnej kontroli pokrycia;
sama ich definicja nie gwarantuje dokładnego poziomu $1-\alpha$.
\end{definicja}

\begin{przyklad}[Dane o wadach: bootstrap może zawieść na brzegu]
\label{ex:bootstrap-zero}
Załóżmy, że w $n=40$ niezależnych próbach Bernoulliego
nie stwierdzono żadnej wady. Wtedy $\widehat p=0$
i rozkład empiryczny skupia całą masę w zerze.
\emph{Każda} próba nieparametrycznego bootstrapu złożona
jest wyłącznie z zer: $\widehat p^*=0$ prawie na pewno,
a oba przedziały z definicji~\ref{def:bootstrap-intervals}
redukują się do $[0,0]$. Taki wynik nie oznacza,
że prawdziwe prawdopodobieństwo wady wynosi zero.
Dla niezależnej próby dwumianowej można zastosować
przedział dokładny Cloppera--Pearsona lub Wilsona
z rozdziału~\ref{ch:confidence-intervals}.
Jest to przykład, w którym zwiększanie $B$ nie naprawi
wad modelu resamplingowego.
\end{przyklad}

\begin{uwaga}[Przybliżenia i diagnostyka]
Bootstrap percentylowy potrafi uwzględniać asymetrię rozkładu
statystyki, ale przy małych próbach, silnym obciążeniu lub na
brzegu przestrzeni parametrów może mieć bardzo złe pokrycie.
Nawet dla regularnych estymatorów nie zawsze ma pokrycie
drugiego rzędu ani nie musi zachowywać granic przestrzeni parametrów.
W praktyce porównujemy co najmniej rozkład replik, liczbę $B$,
stabilność kwantyli i oszacowanie studentyzowane, jeżeli jest dostępne.
Metoda BCa koryguje obciążenie i przyspieszenie; jej pełna teoria
wymaga dodatkowych warunków regularności i nie wynika wyłącznie
z podstawowego CLT bootstrapowego.
\end{uwaga}

\section{Jackknife: usuwanie po jednej obserwacji}
\begin{definicja}[Estymator wariancji jackknife]
\label{def:jackknife}
Niech $n\ge2$, a $T_{(-i)}$ oznacza wartość statystyki
obliczonej z danych bez obserwacji $i$.
Definiujemy $\overline T_{(-)}=n^{-1}\sum_{i=1}^nT_{(-i)}$
oraz
\[
 \widehat V_{\rm jack}=
 \frac{n-1}{n}\sum_{i=1}^n(T_{(-i)}-\overline T_{(-)})^2.
\]
Pierwiastek z $\widehat V_{\rm jack}$ jest błędem standardowym
jackknife; jego poprawność zależy od rodzaju statystyki.
\end{definicja}
\begin{stwierdzenie}[Jackknife dla średniej]
\label{prop:jackknife-mean}
Jeżeli $T_n=\overline X_n$ i $n\ge2$, to
$\widehat V_{\rm jack}=S_n^2/n$.
\end{stwierdzenie}
\begin{proof}
Mamy $T_{(-i)}=(n\overline X_n-X_i)/(n-1)$,
a średnia wartości usuniętych równa się $\overline X_n$.
Wobec tego
\[
 \widehat V_{\rm jack}=
 \frac{n-1}{n}\sum_{i=1}^n
 \frac{(X_i-\overline X_n)^2}{(n-1)^2}
 =\frac{S_n^2}{n}.
\]
\end{proof}

\begin{przyklad}[Dlaczego jackknife zawodzi dla maksimum]
Jeśli $x_{(n)}>x_{(n-1)}$, to spośród $n$ statystyk
maksimum po usunięciu jednej obserwacji $n-1$ wynosi
$x_{(n)}$, a jedna wynosi $x_{(n-1)}$.
Bezpośredni rachunek daje
$\widehat V_{\rm jack}=((n-1)/n)^2
(x_{(n)}-x_{(n-1)})^2$.
Dla niezależnych losowań z $\operatorname{Unif}(0,\theta)$
wariancja maksimum wynosi dokładnie
$\theta^2 n/[(n+1)^2(n+2)]$, ponieważ
$X_{(n)}/\theta$ ma rozkład $\operatorname{Beta}(n,1)$.
Odstęp $D_n=X_{(n)}-X_{(n-1)}$ spełnia
$\Prob(D_n>t)=(1-t/\theta)^n$ dla $0\le t\le\theta$.
Aby to sprawdzić, całkujemy wspólną gęstość
$n(n-1)u^{n-2}/\theta^n$ pary
$(X_{(n-1)},X_{(n)})$ po obszarze
$0<u<v-t<\theta-t$.
Stąd $nD_n/\theta\Rightarrow\operatorname{Exp}(1)$
i wobec tego
$\widehat V_{\rm jack}/\Var(X_{(n)})
\Rightarrow E^2$ dla $E\sim\operatorname{Exp}(1)$.
Granica nie jest stałą jeden, więc jackknife nie daje
zgodnego estymatora wariancji maksimum. Problemem jest
nieregularność tej statystyki, nie niedostateczna liczba
replik.
\end{przyklad}

\section{Testy permutacyjne: dokładność wynika z wymienności}
\begin{twierdzenie}[Dokładny test permutacyjny dwóch grup]
\label{thm:permutation-exact}
Niech $X_1,\ldots,X_m,Y_1,\ldots,Y_k$ będą niezależne
i niech wszystkie mają \emph{ten sam} rozkład $P$,
przy czym $m,k\ge1$, $n=m+k$.
Niech $t$ będzie dowolną mierzalną statystyką przypisania
$m$ z $n$ zaobserwowanych wartości do pierwszej grupy,
a pozostałych $k$ do drugiej. Oznaczmy przez $\mathcal G$
zbiór wszystkich $\binom nm$ przypisań (traktujemy indeksy
jako różne, nawet jeśli wartości się powtarzają).
Dla przypisania obserwowanego $g_0$ definiujemy
\[
 p_{\rm perm}=
 \frac{\#\{g\in\mathcal G:t(g)\ge t(g_0)\}}{\binom nm}.
\]
Wtedy dla każdego $\alpha\in[0,1]$
$\Prob_{H_0}(p_{\rm perm}\le\alpha)\le\alpha$.
Przy braku remisów wartości $t(g)$ i na odpowiedniej
siatce poziomów zachodzi równość.
\end{twierdzenie}
\begin{proof}
Warunkowo względem nieuporządkowanego multizbioru
zaobserwowanych wartości, pod hipotezą wspólnego rozkładu
wszystkie przypisania \emph{indeksowanych obserwacji}
są jednakowo prawdopodobne. Uporządkujmy wartości statystyki
$t(g)$ malejąco, dopuszczając remisy. Dla dowolnego
$r\in\{0,\ldots,\binom nm\}$ liczba przypisań,
dla których liczba wartości $t$ nie mniejszych od obserwowanej
nie przekracza $r$, wynosi co najwyżej $r$.
Dzieląc przez $\binom nm$ uzyskujemy warunkową nierówność
$\Prob(p_{\rm perm}\le\alpha\mid\text{dane nieuporządkowane})
\le\alpha$, a całkując po danych, tezę.
\end{proof}

\begin{uwaga}[Nie jest to automatycznie test równości średnich]
\label{rem:perm-exchangeability}
\emph{Dokładność} powyższego testu wynika z równości
\emph{całych rozkładów} grup, nie z samej równości średnich.
Jeżeli grupy mają różne wariancje lub inny kształt rozkładu,
to permutacja niewystudentyzowanej różnicy średnich może
mieć niewłaściwy poziom przy $H_0:\E X=\E Y$.
W takiej sytuacji można rozważyć statystykę Welcha,
permutację studentyzowaną z odpowiednimi wynikami
asymptotycznymi albo odporny model z właściwymi
błędami standardowymi. Najpierw trzeba ustalić,
\emph{co dokładnie ma być hipotezą zerową}.
\end{uwaga}

\begin{stwierdzenie}[Test znakowych permutacji dla par]
\label{prop:signflip}
Niech $D_1,\ldots,D_n$ będą niezależnymi zmiennymi,
a ich rozkłady będą symetryczne względem zera,
$D_i\overset d=-D_i$, $i=1,\ldots,n$.
Jeśli oznaczymy obserwowane różnice przez $d_i$,
a przez $t(\varepsilon_1d_1,\ldots,\varepsilon_nd_n)$
statystykę dla $\varepsilon_i\in\{-1,1\}$,
to udział spośród $2^n$ znaków dających wartość
co najmniej obserwowaną definiuje poprawną
(w sensie poziomu nie większego od $\alpha$)
wartość p pod wskazaną hipotezą symetrii.
\end{stwierdzenie}
\begin{proof}
Rozkład wektora $D=(D_1,\ldots,D_n)$ jest niezmienniczy
względem grupy wszystkich zmian znaków, ponieważ
każda współrzędna jest symetryczna i współrzędne są niezależne.
Warunkowo względem orbity wektora pod tą skończoną grupą,
zaobserwowany reprezentant jest jednostajny (licząc
wielokrotności wynikające z zer). Argument rangowy
z dowodu twierdzenia~\ref{thm:permutation-exact}
daje nierówność dla wartości p.
\end{proof}

\begin{uwaga}[Monte Carlo zamiast wszystkich permutacji]
Jeżeli $\binom nm$ lub $2^n$ jest zbyt duże,
losujemy $B$ permutacji niezależnie i jednostajnie.
Przy częstości $b$ permutacji o statystyce co najmniej
obserwowanej używa się wartości
$\widehat p=(b+1)/(B+1)$, nie samego $b/B$.
Korekta ta uwzględnia obserwowaną konfigurację
w porównaniu rangowym; w szczególności nigdy nie
raportujemy wartości p równej zeru tylko dlatego,
że przy skończonym $B$ nie trafiliśmy skrajniejszego wyniku.
Dokładność zależy od poprawnego mechanizmu randomizacji.
\end{uwaga}

\section{Zależność w czasie: dlaczego potrzebny jest bootstrap blokowy?}
\label{sec:block-bootstrap}
Załóżmy, że mamy stacjonarny szereg $(X_t)_{t\in\mathbb Z}$
z sumowalnymi autokowariancjami $\gamma_h=\Cov(X_0,X_h)$.
Wówczas, zgodnie z rozdziałem~\ref{ch:time-series},
\[
 n\Var(\overline X_n)\longrightarrow
 \sigma_\infty^2=\gamma_0+2\sum_{h=1}^\infty\gamma_h.
\]
Naiwny bootstrap losujący pojedyncze obserwacje odtwarza
jedynie marginalny rozkład empiryczny. Jeżeli ponadto
szereg jest ergodyczny i $\E X_0^2<\infty$, to
$n\Var^*(\overline X_n^*)=(n-1)S_n^2/n\to\gamma_0$
prawie na pewno, a prawdziwa granica
$n\Var(\overline X_n)\to\sigma_\infty^2$
może być inna. Same założenia stacjonarności i sumowalności
autokowariancji nie zapewniają ergodyczności ani
zbieżności empirycznej wariancji do $\gamma_0$.
Dla modelu AR(1) z korelacją $\rho^h$, $0<\rho<1$,
iloraz poprawnej wariancji asymptotycznej do wariancji
naiwnej wynosi $(1+\rho)/(1-\rho)$.

\begin{definicja}[Bootstrap ruchomych bloków]
Dla próby $x_1,\ldots,x_n$ oraz długości bloku
$1\le\ell\le n$ definiujemy $n-\ell+1$ bloków
$(x_j,\ldots,x_{j+\ell-1})$, $j=1,\ldots,n-\ell+1$.
Losujemy niezależnie i jednostajnie bloki ze zwracaniem,
łączymy je i obcinamy wynik do długości $n$.
Tak otrzymana próba ma długość $n$ i zachowuje
część lokalnych zależności wewnątrz bloków.
\end{definicja}

\begin{uwaga}[Jak dobrać długość bloku?]
W typowych twierdzeniach dla słabo zależnych,
stacjonarnych szeregów wymaga się
$\ell=\ell_n\to\infty$ oraz $\ell_n/n\to0$,
ale nie są to warunki \emph{wystarczające bez dodatkowych
hipotez} mieszania, momentów i właściwego funkcjonału.
Zbyt krótkie bloki zrywają zależność; bardzo długie
zostawiają niewiele efektywnie różnych bloków.
Pomiary powtarzane od jednego pacjenta wymagają z kolei
bootstrapu \emph{klastrowego}, nie blokowego w czasie,
o ile zależność wynika z przypisania do pacjenta.
\end{uwaga}

\section{Dobór metody do problemu}
\begin{center}\small
\begin{tabular}{>{\raggedright\arraybackslash}p{.25\textwidth}>{\raggedright\arraybackslash}p{.29\textwidth}>{\raggedright\arraybackslash}p{.30\textwidth}}
\toprule
\textbf{Sytuacja} & \textbf{Metoda do rozważenia} & \textbf{Warunek krytyczny}\\
\midrule
Niezależne obserwacje, gładka statystyka & Bootstrap nieparametryczny & Reprezentatywność próby; regularność statystyki\\
Silny model parametryczny, mała próba & Bootstrap parametryczny & Model dopasowany poprawnie\\
Średnia i gładki funkcjonał wielu danych & Jackknife lub bootstrap & Rozkład statystyki nie jest skrajnie nieregularny\\
Dwie grupy o jednakowym rozkładzie pod $H_0$ & Test permutacyjny & Wymienność pod hipotezą zerową\\
Różnice w parach, symetria wokół zera & Sign-flip & Niezależność i symetria różnic\\
Obserwacje zależne w czasie & Bootstrap blokowy & Stacjonarność, słaba zależność i długość bloku\\
Dane klastrowe & Bootstrap klastrów & Jednostką resamplingu jest cały klaster\\
\bottomrule
\end{tabular}
\end{center}

\section{Laboratorium: bootstrap, permutacje i zależne obserwacje}
\begin{lstlisting}[language=Python,caption={Kompletny program: bootstrap średniej, dokładne permutacje i bootstrap blokowy (Python, NumPy i SciPy).}]
import itertools
import numpy as np
from scipy.stats import norm

rng = np.random.default_rng(2026)
x = np.array([9., 11., 10., 12., 10., 9., 13., 8.])
B = 20_000
idx = rng.integers(0, len(x), size=(B, len(x)))
boot = x[idx].mean(axis=1)
print("mean:", x.mean(), "bootstrap SE:", boot.std(ddof=1))
print("percentile CI:", np.quantile(boot, [0.025, 0.975]))
print("classical SE:", x.std(ddof=1) / np.sqrt(len(x)))

# Exact permutation test of equal distributions, small samples.
a = np.array([1.0, 2.0, 4.0])
b = np.array([5.0, 6.0, 8.0])
pooled = np.r_[a, b]
observed = abs(a.mean() - b.mean())
total = extreme = 0
for comb in itertools.combinations(range(len(pooled)), len(a)):
    mask = np.zeros(len(pooled), dtype=bool)
    mask[list(comb)] = True
    t = abs(pooled[mask].mean() - pooled[~mask].mean())
    total += 1
    extreme += (t >= observed - 1e-12)
print("exact two-sided permutation p:", extreme / total)

# AR(1): naive iid bootstrap versus moving-block bootstrap.
rho, n, ell = 0.75, 320, 16
noise = rng.standard_normal(n)
y = np.zeros(n)
for i in range(1, n):
    y[i] = rho * y[i-1] + noise[i]
starts = rng.integers(0, n - ell + 1, size=(B, (n+ell-1)//ell))
block_indices = (starts[:, :, None] + np.arange(ell)).reshape(B, -1)[:, :n]
block_means = y[block_indices].mean(axis=1)
iid_means = y[rng.integers(0, n, size=(B, n))].mean(axis=1)
print("iid bootstrap SE:", iid_means.std(ddof=1))
print("block bootstrap SE:", block_means.std(ddof=1))
print("theoretical AR(1) inflation, variance:", (1+rho)/(1-rho))
\end{lstlisting}

\begin{lstlisting}[language=Julia,caption={Kompletny program: bootstrap, dokładny test permutacyjny i resampling bloków (Julia, Random i Statistics).}]
using Random, Statistics
Random.seed!(2026)
x = Float64[9, 11, 10, 12, 10, 9, 13, 8]
B = 10_000
boot = [mean(x[rand(1:length(x), length(x))]) for _ in 1:B]
println("bootstrap mean SE: ", std(boot))
println("percentile CI: ", quantile(boot, [0.025, 0.975]))

# All 3-element subsets of six indices, implemented without packages.
a = Float64[1, 2, 4]
b = Float64[5, 6, 8]
values = vcat(a, b)
obs = abs(mean(a) - mean(b))
extreme = 0
total = 0
for i in 1:4, j in (i+1):5, k in (j+1):6
    chosen = [i, j, k]
    other = [z for z in 1:6 if !(z in chosen)]
    stat = abs(mean(values[chosen]) - mean(values[other]))
    extreme += stat >= obs - 1e-12
    total += 1
end
println("exact permutation p: ", extreme / total)

rho, n, ell = 0.75, 320, 16
y = zeros(n)
noise = randn(n)
for i in 2:n
    y[i] = rho*y[i-1] + noise[i]
end
nb = cld(n, ell)
block = zeros(B)
for b in 1:B
    starts = rand(1:(n-ell+1), nb)
    sample = reduce(vcat, [y[s:(s+ell-1)] for s in starts])[1:n]
    block[b] = mean(sample)
end
println("block bootstrap SE: ", std(block))
\end{lstlisting}

\begin{zadanie}[Próba z zerową liczbą sukcesów]
W $n=20$ niezależnych próbach Bernoulliego nie obserwowano sukcesów.
Opisać przedział percentylowy nieparametrycznego bootstrapu
i wskazać alternatywę o poprawnym skończonopróbkowym pokryciu.
\par\emph{Rozwiązanie.} Rozkład empiryczny jest skupiony w zerze,
zatem każdy rozkład bootstrapowy $\widehat p^*$ jest skupiony
w zerze i przedział percentylowy wynosi $[0,0]$ dla dowolnego $B$.
Przedział dokładny Cloppera--Pearsona ma dolną granicę zero
i górną $1-(\alpha/2)^{1/n}$ w wersji dwustronnej,
ponieważ $\Prob_p(S_n=0)=(1-p)^n$.
Dla jednostronnego górnego przedziału ufności na poziomie
$1-\alpha$ stosujemy granicę $1-\alpha^{1/n}$.
\end{zadanie}


